%% ============================================================================
%% lemma_chi_multiplicative_v2.tex
%%
%% Round-2 follow-up fix for Lemma \ref{lem:chi-multiplicative}.
%%
%% Reason for the update:
%% In the v1 statement (proofs_major_revision.tex, lines 30--47), the lemma
%% required BOTH summands $L_1, L_2$ to be single-slope. The induction step
%% in Theorem 4.1 (M1) iterates the lemma over (\widehat M_1 \oplus \cdots
%% \oplus \widehat M_{m-1}) \oplus \widehat M_m, where the first operand is
%% already multi-slope after the first iteration. The v2 statement covers
%% exactly this case: one operand may be multi-slope, provided all the
%% slopes that appear (across both operands) are pairwise distinct.
%% ============================================================================

\begin{lemma}[Multiplicativity of the WKB characteristic polynomial
across slope-disjoint formal direct sums]\label{lem:chi-multiplicative}
Let $L_A$ and $L_B$ be scalar linear operators on $\mathbb{C}[x]$ in
$\{\opD, \opS, \opT\}$, each with a single irregular singularity at
$x = \infty$. Suppose that the Newton polygon $N(L_A)$ at $\infty$
consists of edges $E_1^A, \dots, E_{m_A}^A$ of distinct slopes
$-p_i^A / q_i^A$ ($i = 1, \dots, m_A$, $\gcd(p_i^A, q_i^A) = 1$,
$p_i^A, q_i^A \ge 1$), and that $N(L_B)$ consists of a \emph{single}
edge of slope $-p_B / q_B$ with $\gcd(p_B, q_B) = 1$, $p_B, q_B \ge 1$.
Assume the slope-disjointness hypothesis
\begin{equation}\label{eq:slope-disjoint}
   \frac{p_B}{q_B} \;\notin\; \Big\{\frac{p_1^A}{q_1^A},\,
   \dots,\, \frac{p_{m_A}^A}{q_{m_A}^A}\Big\}.
\end{equation}
Let $\widehat{M}(L_A) \oplus \widehat{M}(L_B)$ denote the formal direct
sum of the associated $\mathbb{C}((x^{-1}))$-modules, and let $L_{AB}$
be any scalar operator on $\mathbb{C}[x]$ obtained from a cyclic vector
for the direct sum. Then
\[
   \chi_{L_{AB}}(c) \;=\; \chi_{L_A}(c)\, \chi_{L_B}(c),
\]
where $\chi_{L_A}, \chi_{L_B}$ denote the WKB characteristic polynomials
of $L_A, L_B$ (the product, in the case of $L_A$, of its per-edge
characteristic polynomials over $E_1^A, \dots, E_{m_A}^A$).
\end{lemma}

\begin{proof}
The Newton polygon of a formal direct sum of irregular connections
equals the upper convex hull of the union of the Newton polygons of the
summands (\cite[Ch.~III, \S 1]{sabbah-isomon}; or, for the classical
scalar case, \cite[Ch.~II]{robba-malgrange}). Under
\eqref{eq:slope-disjoint}, this upper hull consists of $m_A + 1$ edges,
exactly one for each slope in $\{p_i^A/q_i^A\}_{i=1}^{m_A} \cup
\{p_B/q_B\}$.

We analyse the $\LCinf$-data along each edge $E$ of $N(L_{AB})$. There
are two cases.

\smallskip
\textit{Case 1: $E$ has slope $-p_i^A/q_i^A$ for some
$i \in \{1, \dots, m_A\}$.}
On the lattice $E \cap \mathbb{Z}^2$, only the summand $\widehat M(L_A)$
contributes leading-order ($\LCinf$) terms at slope $p_i^A/q_i^A$;
the summand $\widehat M(L_B)$ has a single slope $p_B/q_B \ne p_i^A/q_i^A$
and therefore contributes only lower-order ($x^{-1}$-higher-valuation)
corrections to the $\LCinf$-data on $E$. Hence
\[
   \LCinf\!\big(a_k^{(AB)}\big)\big|_{E}
   \;=\;
   \LCinf\!\big(a_k^{(A)}\big)\big|_{E_i^A}
   \qquad \text{for } (k, h) \in E \cap \mathbb{Z}^2,
\]
which by \Cref{def:chi} gives the per-edge identity
$\chi^{(i,A)}_{L_{AB}}(c) = \chi^{(i)}_{L_A}(c)$.

\smallskip
\textit{Case 2: $E$ has slope $-p_B/q_B$.}
By the same argument with the roles of $L_A$ and $L_B$ exchanged, the
$\LCinf$-data on $E$ is contributed solely by the (single edge of)
$\widehat M(L_B)$. Hence
$\chi^{(B)}_{L_{AB}}(c) = \chi_{L_B}(c)$.

\smallskip
Multiplying over all edges of $N(L_{AB})$ and using that for a scalar
operator with multiple slope-distinct edges the total WKB characteristic
polynomial is the product of its per-edge characteristic polynomials,
\[
   \chi_{L_{AB}}(c)
   \;=\;
   \Big(\prod_{i=1}^{m_A} \chi^{(i,A)}_{L_{AB}}(c)\Big)\,
   \chi^{(B)}_{L_{AB}}(c)
   \;=\;
   \Big(\prod_{i=1}^{m_A} \chi^{(i)}_{L_A}(c)\Big)\,
   \chi_{L_B}(c)
   \;=\;
   \chi_{L_A}(c)\, \chi_{L_B}(c),
\]
as claimed.
\end{proof}

\begin{remark}[Iteration over arbitrarily many slope-disjoint summands]
\label{rem:chi-mult-iteration}
The single-slope restriction on $L_B$ is essential only in the
inductive step; \emph{either} operand may be multi-slope provided the
slope-disjointness hypothesis~\eqref{eq:slope-disjoint} (interpreted as
disjointness of slope sets) holds. In particular, an obvious induction
on the number of slopes establishes the multi-fold multiplicativity
statement
\[
   \chi_{L_1 \oplus \cdots \oplus L_m}(c)
   \;=\;
   \prod_{\alpha=1}^{m}\, \chi_{L_\alpha}(c),
\]
where each $L_\alpha$ is single-slope of slope $-p_\alpha/q_\alpha$ and
the slopes are pairwise distinct. This is exactly the form invoked in
the proof of \Cref{thm:M1} (Step 3): one applies the lemma to
$(\widehat M_1 \oplus \cdots \oplus \widehat M_{m-1}) \oplus \widehat
M_m$, with $L_A := L_1 \oplus \cdots \oplus L_{m-1}$ a multi-slope
operator (with slope set $\{p_1/q_1, \dots, p_{m-1}/q_{m-1}\}$) and
$L_B := L_m$ single-slope with slope $p_m/q_m$ distinct from those of
$L_A$. The slope-disjointness hypothesis is exactly the multi-slope
hypothesis $p_i/q_i \neq p_j/q_j$ for $i \ne j$ in the parent operator
$L$.
\end{remark}
